% Protocol-aligned benchmark evaluation. FD001 is simulated and uses artificial
% sensor deletion; it is reported as benchmark evidence rather than field validation.
\subsection{Conditional Cross-View Geometry on C-MAPSS}
\label{sec:results_cmapss}

To test the geometric hypothesis under the shared-model protocol, we evaluated PROP-DIRECT and B-STDPROB on C-MAPSS FD001~\citep{saxena2008cmapss}. FD001 contains 100 training and 100 test engine trajectories generated by NASA's Commercial Modular Aero-Propulsion System Simulation, with 14 non-constant sensor channels after preprocessing. It provides a public simulated reference for the analysis, while remaining distinct from field data.

\subsubsection{Shared-model protocol}

We retained the full 14-channel view and formed two reduced views by deleting three channels from each (11 channels per view; the two deletion sets differ by six channels). For each method, a single model was trained on the union of all three masked views. The 100 training engines were split by engine into 80 fitting and 20 validation engines, so views of the same engine remained in one split. The models use the MSE point-prediction head implemented for this benchmark; 128 dropout passes at test time provide predictive samples. We evaluate the final window of each of the 100 test engines and report energy distances on the raw RUL scale.

\subsubsection{Agreement depends on the view pair}

Table~\ref{tab:cmapss_validation} shows a conditional pattern. PROP-DIRECT reduces the distance to the full view by 18.5\% (90\% CI: 7.8--29.3\%) for View 1 and by 28.9\% (18.4--37.9\%) for View 2. For the two reduced views, however, the ordering reverses: PROP-DIRECT has a 36.6\% larger distance than B-STDPROB (90\% CI: $-62.9$ to $-13.1$\%). The shared model therefore aligns predictions when a reduced view is compared with the full sensor complement, but it does not maintain that alignment when two equally sized views differ in sensor identity.

% Table: C-MAPSS benchmark evaluation (shared-model regime, quality-controlled)
\begin{table}[t]
\centering
\begin{threeparttable}
\caption{\textbf{Cross-view energy distances on the C-MAPSS FD001 simulated
turbofan benchmark under the shared-model protocol.} A single model per
architecture is trained on data that mixes all three observation
configurations, mirroring the synthetic protocol. Lower energy distance
indicates closer agreement between dropout-sampled predictive distributions
for the same engine viewed under different sensor configurations; distances
are on the raw RUL scale (cycles). Bootstrap 90\% confidence intervals in
brackets (2,000 replicates, resampling engines). Negative relative reduction
indicates PROP-DIRECT is \emph{worse}.}
\label{tab:cmapss_validation}
\begin{tabular}{lccc}
\toprule
\textbf{View Pair} & \textbf{PROP-DIRECT} & \textbf{B-STDPROB} & \textbf{Relative Reduction} \\
\midrule
Full vs.\ View 1    & 1.85 [1.61, 2.11] & 2.27 [2.04, 2.51] & 18.5\% [7.8\%, 29.3\%] \\
Full vs.\ View 2    & 1.80 [1.56, 2.05] & 2.54 [2.29, 2.79] & 28.9\% [18.4\%, 37.9\%] \\
View 1 vs.\ View 2  & 1.85 [1.57, 2.17] & 1.36 [1.19, 1.52] & $-36.6$\% [$-62.9$\%, $-13.1$\%] \\
\bottomrule
\end{tabular}
\begin{tablenotes}
\footnotesize
\item Single-view test MAE (cycles), same runs: PROP-DIRECT 16.44 (full) /
15.95 (View 1) / 16.25 (View 2); B-STDPROB 16.27 (full) / 15.42 (View 1) /
15.73 (View 2). The near-identical point accuracy makes this geometry
comparison quality-controlled: differences in cross-view distance cannot be
attributed to differences in posterior informativeness.
\end{tablenotes}
\end{threeparttable}
\end{table}


The single-view MAEs are close across methods. PROP-DIRECT obtains 16.44, 15.95, and 16.25 cycles for the full, View 1, and View 2 conditions, compared with 16.27, 15.42, and 15.73 for B-STDPROB. The reversal in the reduced-view pair is consequently not explained by a large point-accuracy collapse. It is consistent with a more specific sensitivity to which sensors are present, although the benchmark has no oracle posterior with which to identify the mechanism.

This shared-model result is stronger evidence than the earlier specialist comparison because the conditioning pathway is trained across views, but it remains evidence from a simulated benchmark with artificial deletion. It supports a conditional extension of the synthetic finding and sharpens the importance-awareness boundary; it does not establish performance under naturally occurring sensor loss. The archived protocol and outputs make a future study with domain-informed critical channels and field data possible.
